%!TEX root = ../main.tex

% ============================================================
%  CAPÍTULO 4: OBJETIVOS
% ============================================================

\chapter{Objetivos}\label{ch:objetivos}

% ──────────────────────────────────────────────────────────────
\section{Objetivo general}
% ──────────────────────────────────────────────────────────────

Diseñar e implementar un modelo de \textit{nowcasting} de la intención de
voto en elecciones locales colombianas que sea elección-agnóstico, no dependa
de encuestas, use únicamente interacciones en redes sociales de los perfiles
oficiales de los candidatos, y pueda replicarse con herramientas comerciales
accesibles.

% ──────────────────────────────────────────────────────────────
\section{Objetivos específicos}
% ──────────────────────────────────────────────────────────────

\begin{enumerate}
  \item \textbf{OE1.} Replicar y evaluar empíricamente el modelo SOMEN-DC en
        la segunda vuelta presidencial de Bolivia (octubre de 2025).
  \item \textbf{OE2.} Modelar el crecimiento temporal de las interacciones en
        redes sociales usando varios modelos de forma comparativa.
  \item \textbf{OE3.} Construir un \textit{dataset} de entrenamiento
        elección-agnóstico con datos de las elecciones locales colombianas
        de 2023.
  \item \textbf{OE4.} Entrenar y evaluar modelos de clasificación y regresión
        del resultado electoral.
  \item \textbf{OE5.} Documentar las condiciones de viabilidad metodológica
        del modelo propuesto y documentar sus limitaciones.
\end{enumerate}
